\scanpage{7}
\[
\begin{aligned}
H^i({}_pT)&\longrightarrow {}_pH^i(T)\longrightarrow 0\\
0&\longrightarrow H^i(T)_p\longrightarrow H^{i+1}({}_pT)
\end{aligned}
\]
on voit que les conditions sont nécessaires.

Pour la suffisance, on tire d'abord le fait que $H^{\deleted{\unsure{n}}i}(G)_{\deleted{\unsure{p}} \unsure{p}}=$ pour $i\geq n+2$, reste à vérifier $H^{n+1}(G)=0$ donc que $H^n(T)\deleted{{}_{\unsure{(p)}}}\longrightarrow H^n(T')\deleted{{}_{\unsure{(p)}}}$ est surjectif. \deleted{\unsure{avec|Bien} \unsure{T} \unsure{De fait. Faisons} Le conoyau est un \unsure{groupe} de $p$-torsion, donc on peut se borner à regarder $H^n(T)_{(p)}\longrightarrow H^n(T')_{(p)}$. C'est une \unsure{bonne|horreur} que $G$ soit $p$-divisible,}

Noter que si $k'/k$ est une extension finie séparable de $k$, alors pour tt $G'$ sur $k'$, on a
\[
H^i(k,\prod_{k'/k}G')\simeq H^i(k,G),
\]
ce qui prouve que l'hypothèse faite sur les tores sur $k$ est \uline{stable} par extension de base séparable. Soit $k'$ une telle extension de degré premier à $p$, alors l'hom trace $\prod_{k'/k}G_{k'}\longrightarrow G$ est \deleted{surjectif} \inserted{un épim.}
% archive sheet number: 6

